\refstepcounter{section}
\section*{Lecture 26: Further Scaling}
\addcontentsline{toc}{section}{Lecture 26: Further Scaling}
The scaling forms used in the previous lectures all come by assuming that the thermodynamic quantities are generalised homogenous functions. So let us discuss a bit about homogenous functions in general. A function $f:\re \to \re$ is called \textit{homogenous} if $$f(\lambda x) = g(\lambda) f(x)$$ 
where $g(\lambda)$ is some function. A simple example is, $f(x) = x^2$ since $f(\lambda x) = \lambda^2 x^2 = \lambda^2  f(x)$, with $g(\lambda) = \lambda^2$. Note that, $g(\lambda)$ cannot be any arbitrary function, it must satisfy some properties. For this, note: 
$$f(\lambda (\mu x)) = g(\lambda) g(\mu)f(x) = f((\lambda \mu)x) = g(\lambda \mu) f(x)$$
which imposes the condition $g(x) g(y) = g(xy)\ \forall \ x,y$. This eventually forces the form $g(x) = x^\rho$ which has an elegant proof, considering $g$ to be differentiable. 
\begin{proposition}
    If $g: \re\to \re$ is a \textit{differentiable} function and $g(x) g(y) = g(xy)\ \forall \ x,y$ then for all $x$, either $g(x)=0$ or $g(x) = x^\lambda$ for some $\lambda$
\end{proposition}
\textit{Proof.} Let any arbitrary $y=y_0$. Differentiating both sides, we get:
$$g'(x) g(y_0) = y_0 g'(xy_0)\ \forall \ x\quad \xlongrightarrow{x=1} \quad  \frac{g'(y_0)}{g(y_0)} = \frac{g'(1)}{y_0}\quad\xlongrightarrow{\text{Integrating}} \quad g(y_0) = A y_0^\lambda$$
where, after differentiation we have taken the $g$ as a function of $y_0$ since $y_0$ was arbitrary and then replaced $g'(1) = \lambda$. The initial condition $g(1) = (g(1))^2 \implies g(1)=0$ or $g(1)=1$. If $g(1) = 0$ then this implies $A=0$, implying that $g$ is identically zero. Otherwise, if $g(1) = 1$ then $A =1$ which gives us $g(x) = x^\lambda$\\[0.2cm]
We can extend the definition of homogenity to arbitrary number of variables. 
$$f(\lambda x_1, \lambda x_2\ldots \lambda x_n)  = g(\lambda) f(x_1,x_2\ldots x_n) \implies f(\lambda \vb{x}) = \lambda^p f(\vb{x})\quad \vb{x} \in \re^{n}$$
The cool thing with homogenous functions is that if we know the value of the function at one point, then we know it at every other point, since given $f(x_0)=c$, we can always choose $\lambda = \frac{y}{x_0}$ such that $f(y) =\qty(\frac{y}{x_0})^\lambda f(x_0) $. We now extend the definition of a homogenous function to generalised homogenous function (GHF); a function $f:\re^n \to \re^m$ is called a GHF if, 
$$f(\lambda^{a_1} x_1,\lambda^{a_2} x_2  ,\ldots, \lambda^{a_n}x_n) = \lambda^k f(x_1, x_2 ,\ldots, x_n)$$
where $\{a_i\}$ and $k$ are some arbitrary numbers. Note that these are not independent, since if we choose an arbitrary $\lambda = \mu^{1/k}$ then we have, 
$$f(\mu^{a_1/k} x_1,\mu^{a_2/k} x_2  ,\ldots, \mu^{a_n/k}x_n) = \mu f(x_1, x_2 ,\ldots, x_n)$$
Then all $a_i \mapsto a_i/k$ which is just a rescaling. The GHF can thus be defined by always choosing $k=1$ and hence the final definition,
$$f(\lambda^{a_1} x_1,\lambda^{a_2} x_2  ,\ldots, \lambda^{a_n}x_n) = \lambda f(x_1, x_2 ,\ldots, x_n)$$
Let us now focus only on the two variable case, which is $F(t,h)$. The static scaling hypothesis states that the thermodynamic potentials are GHFs near the critical point $T_c$. Then the free energy can be written as, 
$$F(\lambda^{a\tund{t}}t, \lambda^{a\tund{h}}h) = \lambda F(t,h)\quad \longrightarrow \quad -\pdv{h}\qty[F(\lambda^{a\tund{t}}t, \lambda^{a\tund{h}}h)] = -\lambda \pdv{h}\qty[F(t,h)] \implies \lambda^{a\tund{h}-1} m(\lambda^{a\tund{t}}t, \lambda^{a\tund{h}}h) =m(t,h)$$
Now, suppose we take $T>T_c$. Then we choose $\lambda$ such that $\lambda^{a\tund{t}} = -\frac{1}{t}\implies \lambda = (-t)^{-1/a\tund{t}}$. Taking the zero-field limit we get: 
$$(-t)^{\beta} \sim m(t,0) = (-t)^{-(a\tund{h}-1)/a\tund{t}}m(1,0) \implies \boxed{\beta = \frac{1-a\tund{h}}{a\tund{t}}}$$ 
Now, suppose we take $T=T_c $ and choose $\lambda$ such that $\lambda^{a\tund{h}}=\frac{1}{h}\implies \lambda = h^{-1/a\tund{h}}$. Substituting this in the scaling form above, we get:
$$h^{1/\delta} \sim m(0,h) = h^{-(a\tund{h}-1)/a\tund{h}}m(0,1)\implies \boxed{\delta = \frac{a\tund{h}}{1-a\tund{h}}}$$
Solving these we get,
\begin{equation}
    a\tund{h} = \frac{\delta}{1+\delta} \qquad a\tund{t} = \frac{1}{\beta(1+\delta)} 
\end{equation}
We get $a\tund{h}$ and $a\tund{t}$ in terms of the critical exponents (and vice versa). Now, differentiating $m$ with respect to $h$ gives us the susceptibility:
\begin{align*}
   \lambda^{2a\tund{h}-1} \chi(\lambda^{a\tund{t}}t, \lambda^{a\tund{h}}h) = \chi(t,h)
\end{align*}
For $T>T_c$, taking $h=0$ and choosing $\lambda$ such that $\lambda^{a\tund{t}}=\frac{1}{t}$, we get the following
$$t^{-\gamma'}\chi(t,0) = t^{-(2a\tund{h}-1)/a\tund{t}} \chi(1,0) $$
from which we can read off the exponent as,
$$\gamma' =  \frac{(2a\tund{h}-1)}{a\tund{t}}= \frac{\frac{2\delta}{1+\delta}-1}{\frac{1}{\beta(1+\delta)}} = (\delta-1)\beta \implies \beta \delta = \beta+\gamma$$
which was the Widom's scaling law that we had earlier derived.  One thing to note here is that, if we had taken $T<T_c$ then also the same calculation follows, which would give us the same expression for $\gamma$. Hence we can conclude that $\gamma=\gamma'$

Doing the same thing for the specific heat will lead to the Rushbrooke relation (which I skip since it has become very boring at this point of time!).
\subsection{A small note}
The scaling laws that we derived manifest as inequalities when derived from thermodynamic principles, considering stability of physical systems. Let us try to derive the Rushbrook (in)equality! 

Consider the first law of thermodynamics for a magnetic system with field $H$ and and magnetisation $M$:
\begin{equation}
    \dd U = T\dd S - M \dd H \implies T = \qty(\pdv{U}{S})_{H} \quad M = -\qty(\pdv{U}{H})_{S}
\end{equation}
We can now define the specific heat at constant magnetisation and at constant field as, 
\begin{align*}
    C_H = T \qty(\pdv{S}{T})_H\qquad   C_M = T \qty(\pdv{S}{T})_M
\end{align*}